% !PRINTED 679
\weberpage{679}
Bezeichnen wir nämlich mit $x_i'$ ganze rationale Funktionen von $z'$, höchstens vom Grade $m$, so ist
\weq{7}{x_i = z^m x_i'}
eine ganze rationale Funktion von $z$. Wäre nun
\[
\lambda_s' = \sum_{0,\,s-1}^{i} x_i'\lambda_i',
\]
so würde aus (6) und (7) folgen:
\[
z^m\lambda_s = \sum_{0,\,s-1}^{i} x_i z^{r_s-r_i}\lambda_i,
\]
und dies ist unmöglich, weil die $\lambda_i$ linear unabhängig sind; und es kann auch keine Funktion $\lambda_s'$ von niedrigerem Exponenten als $r_s$ geben, weil man daraus eine Funktion $\lambda_s$ von niedrigerem Exponenten herleiten könnte.

Die Körperdiskriminanten $D_z, D_{z'}$ in bezug auf die Variable $z$ und $z'$ sind
\[
\begin{aligned}
D_z &= \Delta(\lambda_1, \lambda_2, \ldots, \lambda_n), \\
D_{z'} &= \Delta(\lambda_1', \lambda_2', \ldots, \lambda_n'),
\end{aligned}
\]
und daraus folgt nach (6):
\weq{8}{D_{z'} = z'^{2(r_1+r_2+\cdots+r_n)}D_z.}

Ist also $\delta$ der Grad von $D_z$, so ist
\[
2(r_1+r_2+\cdots+r_n) - \delta
\]
die Anzahl der verschwindenden Wurzeln von $D_{z'}$, und es ergibt sich nach \S\,190, 5.:
\weq{9}{w_z = 2(r_1+r_2+\cdots+r_n).}

7. \gesperrt{Die Verzweigungszahl $w_z$ ist also immer eine gerade Zahl.}

\begin{center}
\textbf{\S\,194. Differentialquotienten.}
\end{center}

Die Differentialquotienten können wir, wo wir keinen Gebrauch von der \gesperrt{Stetigkeit} machen, nicht in der gewöhnlichen Weise einführen. Wir nehmen zwei Funktionen $\alpha, \beta$ aus $\Omega$, zwischen denen die irreduzible Gleichung
\weq{1}{F(\alpha, \beta) = 0}
besteht, bezeichnen mit $F'(\alpha), F'(\beta)$ die abgeleiteten Funktionen und \gesperrt{definieren} den Differentialquotienten als Funktion des Körpers $\Omega$ durch die Formel:
\weq{2}{\left(\frac{d\alpha}{d\beta}\right) = -\frac{F'(\beta)}{F'(\alpha)}.}
%CONT>>
